\section{Esoteric Paganism}

The revelation given to the rishis or seers of ancient India was that the empirical ego is not the real self. Of course, the pagan revivalists ignore this and content themselves with putting up statues of Hindu gods, rather than engaging in the difficult process of detaching from the empirical self. For men today, that is the opposite of salvation, for what they want is for the empirical ego to be made more comfortable and to achieve satisfaction. Because of technology, the illusion that such an outcome may be permanently possible takes root. Many illnesses can be cured or at least treated effectively. The effects of old age are diminished because the body has not undergone a lifetime of hard labour in most cases and chronic conditions can be managed. Finally, Death, which came at random among all age groups, is relegated more and more to the end of a longish life. Thus men today keep sickness, old age, and death out of their conscious thoughts for the most part.

Buddha, following and developing the Vedic tradition, saw things in their full impact. \textbf{Valentin Tomberg} wrote:

\begin{quotex}
A Buddha is one who is fully awakened. ... All human beings know about the facts of birth, sickness, old age, and death, and yet live as if they knew nothing of them. But Gautama Buddha became completely awake to these facts, awake to their range and significance right down to their ultimate consequences. ... The Buddhas of the past were fully awake to the reality of the Fall into sin, and to the sorrows of mankind and of the world which is sunken in error, suffering, and death.
\end{quotex}


So Tomberg points to three great revelations. The ideal of the pagan world was the Wise Man, Sage, or the Philosopher, as the one who remembers, i.e., the one who overcomes forgetting. The Buddhas are fully Awakened Ones who overcome sleep. Yet the noble path revealed by the Buddha could not overcome death. So that brings us to the Christian revelation which we will try to understand in its inner, or esoteric, nature.

\paragraph{Philosophy in the Hellenic World}


Philosophy used to be regarded as a way of life leading to Wisdom; no professor of philosophy today can be regarded as ``wise'' in any sense. Plato, Aristotle, and the Stoics were all considered to have parts of the true philosophy. This is not unlike the six philosophical schools of Hinduism, just that, unfortunately, it was not formulated so precisely.

``Sham'' philosophy was the pseudo-philosophy of the sophists, epicureans, materialists, atheists, etc. When St. Paul criticizes philosophy or the traditions of men (Colossians 2:8\footnote{\url{http://www.biblegateway.com/passage/?search=Colossians+2\%3A8\&version=DRA}}), that is what he had in mind.

Academics are fond of saying that the early Christians ``used'' Greek philosophy to express their teaching as it was the most sophisticated thought available to them. That gives them permission to then force Christian teaching onto all sorts of ludicrous ideologies, or better said, ``sham philosophies.'' Of course, the real situation is quite different. They saw pagan Hellenic philosophy as containing essential truths. It was never part of Christian teaching to create a new metaphysical system. \textbf{Vladimir Solovyov} explains:

\begin{quotex}
The great fathers of the ancient Church affirmed that even before Christ's incarnation, the same divine Reason that was revealed in Him enlightened with the eternal truth the inspired wise men of paganism, these \emph{Christians before Christ}.

Although the close inner connection between Alexandrian theosophy and the Christin doctrine is one of the firmly established theses of Western scholarship, for one reason or another, this perfectly correct thesis does not enjoy common acknowledgement in our theological literature.
\end{quotex}


Solovyov promised to write about the revelations of \textbf{Thoth} or \textbf{Hermes Trismegistus} in its relation to Christian doctrines.

Unfortunately, he never got around to it, but that task has since been accomplished by Valentin Tomberg in his meditations.

\paragraph{Exoteric and Esoteric Philosophy}


The Church Fathers continued to call their own teachings ``philosophy'', although they emphasized that true philosophy must be ``lived''. What the means can be found in the \emph{Philokalia}. Because of this they drew the following distinction:

\begin{itemize}


\item Exoteric philosophy comprised ancient Greek philosophy

\item Esoteric philosophy is identical with the Christian religion
\end{itemize}


So from their point of view the Christian religion was entirely esoteric, and was actually the inner aspect of Hellenic philosophy as exoteric. With this understanding the \emph{Philokalia} and Dante become much clearer. We have alluded to Augustine's insight that the external Logos of the pagans was the same as the inner Logos, born again in the soul, of the Christians.

This has been lost, as Christianity became more exoteric. The path to Wisdom, or the Love of God, was grounded in three factors:

\begin{enumerate}


\item Divine revelation, keeping in mind that there were four levels of interpretation, from the literal to the esoteric

\item Sacred Tradition

\item Inner experience, which lifts one above the senses and above reason
\end{enumerate}


Without point (3), there is only an exoteric teaching. More has been lost over time. The Roman Church was based on (1) and (2), although there is a continuous history of saints and mystics who understood (3); their writings should be read often. The Reformation insisted on point (1) alone, and even worse, accepted only the literal level of interpretation.

For references on this topic, please see \textit{The Hellenic-Christian Philosophical Tradition}\footnote{\url{http://ibmgs.org/ancientphil.html}}, by \textbf{Constantine Cavarnos}.


\flrightit{Posted on 2014-03-03 by Cologero}

\begin{center}* * *\end{center}

\begin{footnotesize}
\begin{sffamily}
\texttt{Constantine Aetos on 2014-03-03 at 13:54 said:}

``This has been lost, as Christianity became more exoteric. The path to Wisdom, or the Love of God, was grounded in three factors:''

And that is why Eastern Orthodoxy is the only column left in Christendom.

\hfill

\texttt{Michael on 2014-03-03 at 18:25 said:}

Cologero, do you think that Guenon drew too sharp a distinction between initiation and mysticism? In Christianity is true mysticism and initiation one and the same?


\hfill

\texttt{Paulo on 2014-03-03 at 23:36 said:}

``Yet the noble path revealed by the Buddha could not overcome death.''

You say this because the Buddha could not explain death,could not fill the void?

\hfill

\texttt{Cassiodorus on 2014-03-04 at 18:12 said:}

Hi Cologero,

You have always been gracious in granting me a certain degree of latitude for ``chitchat'' as you are wont to call it. A few years ago, I became familiar with the writings of the Perennialist school which resonated me with me profoundly. I then slowly migrated back to the religion of my childhood, Roman Catholicism. But as I became more immersed in the teachings of the Catholic faith, the more difficult I have found it to affirm both the universalism of Perennialism and the central doctrines of Christianity. For me, the question of whether or not Christianity reveals something\emph{``more''} than the other great Revelations is extremely important. If it does not, then I I'm hard pressed to reconcile the Incarnation and Trinitarian theology with Perennialist metaphysics which regards the Trinity on the level of being- not beyond being (Nirguna Brahman) and Jesus Christ as just one of multiple ``avataric descents''.

In response, Traditionalists usually say that they are certainly not trying to convert anyone who is at peace with God- exoterism, after all, has its rights. But, there \emph{is} a transcendent unity of religions- esoterism transcends theology from \emph{within}, not from without. No everyone is a jnani, nor should we expect them to be.

Traditional Catholics generally react to this pov by charging Perennialism with being a modern mutant form of gnosticism. Their reason is that the Perennial Philosophy regards ignorance as the great sin- ignorance of the supreme Truth that duality is an illusion and therefore that the heart of spirituality is realization. Traditional Christiainty rejects this (says my Catholic interlocutors) on the grounds that only the saving work of the Incarnate Christ reconciles us to God. Moreover, under no circumstance can the Christian accept moksha, \emph{identity} with the Divine. Christian mysticism, union with Christ , transcends other Traditions because the Logos, the Second Person of the Trinity, \emph{became man}. Perennialism is basically a failed attempt to ``squeeze'' classical theism into a nondual framework which fundamentally changes the significance of Jesus Christ.

For what it's worth, I hesitated to post this, especially since I've been beating this horse for a while- but I just can't seem to make peace with it. I am not trying to be argumentative and I'm certainly not trying to rile anyone up. Any insight on this matter would be appreciated.

Vincit Omnia Veritas

\hfill

\texttt{Frater M on 2014-03-05 at 00:59 said:}

Guenon expressed in a brief footnote in his ``Perspectives'' that it seemed likely that St. Ignatius' ``Spiritual Exercises'' were an ``ascetic'' and ``active'' way, likely inspired (in part) by Islamic Esoterism. Others since Guenon, have also explored this idea (independent of Guenon), and arrived at similar conclusions.

Despite any of this, ``The Spiritual Exercises'', imo, can stand on their own, and speak for themselves. St. Ignatius describes them as ``way'', in the sense of pilgrimage, with the pilgrimage being ones' life, and not simply ending after the forty-day retreat, but continuing ever onward, where we so effectively learn the ``exercises'' that eventually, our entire lives constitute the ``retreat'' (may we all obtain the grace to reach such a point).

A goal, or perhaps better stated, a grace possibly afforded the spiritual athlete of these exercises, is ``revelation''--coming to knowledge of ourselves, the meaning of our individual lives, and what needs to be further accomplished in our lives, to allow God to fully self-disclose Himself through us.

Perhaps my own explanation here of the ``Exercises'' is far inadequate, and incomplete; however, the process of the exercises, in general, seem to my mind, to satisfy most of what people are after when seeking an ``esoterism''--and, they are something quite common, in plain sight--yet often glossed over. They are also a very living, and ever growing tradition, within the bosom of Western Christianity/Roman Catholicism.

\hfill

\texttt{Cologero on 2014-03-05 at 07:14 said:}

Yes, Cassiodorus, perhaps it may be more helpful for you to interlocute with what has been written here rather than those anonymous third parties. If you become reconciled to God, is that a real change in your being or merely an imputed reconciliation? If the former, then no interlocution is going to help you. Is reconciliation some impersonal mechanical or juridical process brought about by Christ passively, or is something more involved?

Clearly, we have been relying on Tomberg and Solovyov to properly relate the various traditions. Although we have not said it all yet, don't you have anything to say about the actual text written her beyond the chitchat that is now occupying your mind. Remember the First Trial.

The problem with ``perennialism'' is that it is ahistorical. It teaches the doctrine of cycles, but avoids it in practice by attempting to transcend the world process. But the world process is not an ``illusion'' as you apparently believe, but rather a reflection of the divine, even if imperfect. We have emphasized this over and over both here and on the Tarot blog. There is also a reason for the imperfection.

So, no, unlike the Eastern teachings that want to annihilate the self, something certainly possible and it is the limit of their teaching for the reasons given by Tomberg. Nevertheless, the revelation in the Vedanta of the higher self is valid and important. So, yes, realization trumps mere chitchat.

Beyond that, our goal is regeneration, i.e., the return to the state of consciousness prior to the Fall; that is the way of Christ, but you must follow that path. That is not self-annihilation, but self-affirmation. Oddly enough, Evola and Serrano seem to understand that, which is why we have bothered with them.

\hfill

\texttt{Avery Morrow on 2014-03-05 at 09:36 said:}

Cassiodorus, it is correct to recognize that Guenon has no epistemology. He awakens an intuitive sense that traditional forms serve some higher purpose, but he provides no bounds by which to distinguish these forms perfectly from their mutations, nor does he resolve ancient teachings such as the uniqueness of Christ. Now Tomberg and Solovyov are being employed in this post to provide us with some of those limits. It is worth discussing with your friends whether the points being raised by these writers are discernibly orthodox to them.

\hfill

\texttt{Paulo on 2014-03-06 at 07:29 said:}

Why do you say that the buddha could not overcome death?

\hfill

\texttt{Cassiodorus on 2014-03-07 at 10:41 said:}

Cologero and Avery,

Thank you very much for your comments and insights. I have much to contemplate during this Lenten season.

Kind Regards

\hfill

\texttt{Cologero on 2014-03-10 at 14:45 said:}

First of all, Exit, the idea of ``types'' in Scripture has been known from the beginning. Tomberg has developed the idea. This begs the question of where the ``types'' came from in the first place. If the types are also human constructs, then they would not remain constant over the centuries. So actually archetypes and prototypes are revelations from God regarding the inner structure of the cosmos. Human history follows certain patterns based on such a typology.

If you would seriously like to study the Jewish teaching on the end of the world, check out ``The Way of God'' by Luzzatto, particularly the section ``Israel \& The Nations''. Luzzatto said that the gentile nations are on a lower level than Israel because they rejected the Torah. There is no need for any elaborate plot.

\hfill

\texttt{Frater M on 2014-03-11 at 01:49 said:}

This is really off-topic; however, Algis Uzdavinys didn't really ``write'' that Hyksos material mentioned by ``Exit'' (found in ``Philosophy as a Rite of Rebirth, pp 30-33)--he was quoting Jan Assman, and it isn't really much more than a blurb in the book, or all of Uzdaviny's writings. I mention this because ``Exit's'' comment makes Uzdaviny's appear as a dry academic, ``historian of religion'', who likely fails to grasp any higher meanings. Such is not the case--in fact, Uzdaviny's revisits aspects of the Hebrew tradition, favorably, in his later book ``Ascent to Heaven''.

\hfill

\texttt{Cologero on 2014-03-17 at 00:31 said:}

Without endorsing these particular figures or points of view, the idea of Christianity as esoteric paganism arises in other streams, casting doubt on the notion that Wotanism represents the primal Nordic tradition. First of all, we have pointed out several times that Herman Wirth's researches showed that monotheism was at the source of the Nordic tradition.

Both Karl Wiligut and Guido von List also believed that Wotanism was the exoteric expression of the esoteric ``Irmin-Christianity''.

\hfill

\texttt{burningfireeye on 2021-03-03 at 16:06 said:}

Angeology cannot be compared with paganism. Do we sacrifice, do we make offerings to an idol of St. Michael the Archangel? Did Hinduism consist of sacrificing to pagan idols from its very inception or is it a degeneration? If it is not a degeneration, how can it represent the primordial tradition?

\hfill

\texttt{Cologero on 2021-03-03 at 22:26 said:}

burning fire has apparently affected your eyesight adversely. Where, oh where, did the topic of angelology come up? Or sacrifice? Or Michael? Where exactly is there a comparison?

They say ignorance is bliss, but most of the time it is just the result of laziness. 

Please read carefully \textit{Treatise on Separate Substances by Saint Thomas Aquinas}\footnote{\url{https://www.gornahoor.net/?page_id=13278}}, in which he does indeed discuss angels and paganism.

\end{sffamily}
\end{footnotesize}

